\section{Corrección de pendientes tras la auditoría de bloques I--IV}
\label{sec:correcciones_pendientes}
La auditoría reprodujo dos debilidades que las 140 pruebas originales no
detectaban: aceptar una traducción equivocada como explícita y aprobar con un
veredicto antiguo sin citas. Se conservan las capturas históricas; las nuevas
pruebas y resultados se guardan en \code{validation/pending_fixes/}.

\subsection{P07--P08/E13: correspondencia de restricciones y abstención}
Antes bastaba que una cita existiera y contuviera una negación: ``sin romance''
podía vincularse al descriptor \emph{violence}. Ahora \code{intent.py} exige
igualdad de frase normalizada o una equivalencia explícitamente revisada. Las
equivalencias son específicas de modalidad y no se aceptan porque el modelo
declare que son correctas. No se eliminan acentos para comparar evidencia.
Una traducción desconocida, ampliación no justificada o inversión de polaridad
provoca error de validación y aclaración mediante el fallback seguro. También
se revalidan restricciones explícitas heredadas de un plan previo.

Se exige que una cita negativa incluya su negación al comienzo y conserve todos
los conceptos de una conjunción. Una cita corta seguida de continuación de lista
no acredita el alcance completo. Se rechazan dobles negaciones, construcciones
``no solo'' y contradicciones entre restricciones explícitas. Los casos
acústicos completamente reconocidos mantienen su gramática determinista.
El prompt pide descriptores mínimos y aclaración cuando no pueda verificarse
una traducción. El control de alcance no pretende resolver toda la gramática.

\textbf{Por qué es mejor.} Se elimina el camino demostrado de etiquetar una
traducción incorrecta como evidencia verificada. La comprobación local no depende
de que otro modelo confirme al primero. \textbf{Coste y alternativa.} Se pierde
cobertura: paráfrasis válidas no incluidas pueden pedir aclaración. Es preferible
a buscar ocultando una restricción; un verificador semántico más amplio sólo debe
añadirse con corpus anotado, tasas de error y abstención medidos. E13 queda
corregido para este contrato conservador; la fidelidad general sigue pendiente
de evaluación humana P04/P10/P12.

\subsection{P10: evidencia uniforme para decisiones del juez}
El cliente ya no registra herramientas de exclusión. Sólo se aplican decisiones
JSON estructuradas con motivo y citas verificables dentro del texto suministrado.
Tanto aprobar como rechazar requieren evidencia no vacía, incluso si la petición
contiene únicamente negativos. Una cita prueba presencia textual, no la
corrección semántica ni por sí sola la ausencia de un tema en toda la canción.
La cobertura incompleta sigue impidiendo aprobar la letra completa.

Un texto antiguo o una llamada de herramienta provoca un único reintento de
formato; si persiste, el juez se abstiene. Si las citas faltan o son inventadas
en ambos intentos, también se abstiene. Los fallos de transporte y otras salidas
malformadas conservan estado de error, sin multiplicar llamadas al proveedor.
El orquestador convierte además una decisión marcada \texttt{legacy} en abstención:
no puede convertirse en aprobación seleccionada ni rechazo semántico.

Las pruebas antiguas de ranking se actualizan para emitir decisiones JSON con
citas; las de compatibilidad comprueban ahora abstención, aislamiento y ausencia
de mutación. Las nuevas regresiones verifican recuperación tras reintento,
negativos sin evidencia, herramientas antiguas, traducción incorrecta, negaciones
compuestas y que una interpretación incierta no llegue a la búsqueda.
El coste posible es más abstenciones y una segunda llamada en respuestas antiguas;
el máximo sigue siendo dos por candidato y el refinamiento continúa desactivado
por defecto. Las citas no sustituyen la calibración humana del juicio.

\subsection{P04/P10/P12: evaluación humana preparada, sin etiquetas inventadas}
\code{evaluation/annotations.py} exporta consultas, metadatos y letras del snapshot
con grados y violaciones vacíos. Para audio exige escucha: no permite inferir
relevancia sonora a partir de la letra. La importación exige declaración de
revisión humana, anotador, consultas idénticas y generación/hashes coincidentes;
valida IDs, grados y cobertura. Esa declaración es atribución, no prueba automática
de que una persona revisó correctamente. Los archivos existentes no se sobrescriben.
\code{evaluation.run --qrels} permite usar la revisión sin modificar el archivo
original; una generación distinta bloquea la evaluación.
Se conserva la distinción entre desconocido y grado cero. No se fabricaron
anotaciones, por lo que relevancia y calibración permanecen abiertas.

\subsection{P16: memoria del proceso y perfil real ejecutable}
El perfil de control ahora registra working set RSS mediante la API de Windows
(o \code{/proc/self/statm} en Linux), muestras de memoria y pico del proceso
durante su vida. El pico de muestras puede perder aumentos breves; el pico de
vida incluye inicialización. Los contadores CUDA, si están disponibles, cubren
asignaciones Torch, no toda la memoria del controlador. No se importa Torch
exclusivamente para medir un arranque bloqueado. No se añadieron dependencias.

\code{evaluation/real_profile.py} ejecuta el runner real con muestreo de RSS, sin
sustituir modelos por controles. Los pesos sólo se buscan localmente por defecto;
la descarga requiere la opción explícita \code{--allow-model-downloads}.
La ejecución realizada quedó bloqueada por \code{GROQ_API_KEY} ausente. Su memoria
es de arranque, no de inferencia. El perfil de control sí mide RSS real de ese
proceso, pero no permite extrapolar uso de memoria o latencia de modelos reales.

\subsection{Estado de cierre y prioridad restante}
Los dos controles de la auditoría dejan de reproducir los fallos. La suite y las
series simuladas se vuelven a capturar con hashes de fuentes actuales, conservando
los informes de bloques anteriores. La documentación TeX se comprueba
estáticamente; no se declara compilación PDF ni revisión visual nueva de la web.
La condición pendiente es configurar el proveedor localmente y realizar la
anotación humana, después ejecutar y comparar modelos reales antes de activar
refinamiento o afirmar calidad/rendimiento. E17--E19 conservan su prioridad y
ubicación del bloque V; esta corrección no realiza ingesta real ni limpieza de datos.
